% =====================================================================
% Chapter 4: Results
%
% Covers:
%   4.1 Reporting Conventions
%   4.2 Baseline Establishment and Speaker-Disjoint Protocol Validation
%   4.3 Exploratory Model Capacity Comparison
%   4.4 Research Question 1: Curriculum Learning
%   4.5 Research Question 2: Data Augmentation
%   4.6 Out-of-Domain Evaluation (Google FLEURS)
%   4.7 Computational Cost Analysis
%   4.8 Integrated Summary of Findings
%   4.9 Limitations
% =====================================================================

\graphicspath{{images/}}

\begin{document}

\section{Results}

\subsection{Introduction}

This chapter presents the empirical findings of the experiments described in the
Methodology chapter. All results are reported on the primary evaluation metrics
defined in Section~3.8: Word Error Rate (WER), which measures the proportion of
incorrectly recognised words relative to the reference transcription, and
Character Error Rate (CER), which applies the same edit-distance calculation at
the character level. Lower values indicate superior recognition performance for
both metrics. Where relevant, training runtime (wall-clock seconds) and peak
Graphics Processing Unit (GPU) memory consumption are reported to characterise
computational cost.

The results are organised by research question. For each experimental condition,
three independent training runs were conducted using random seeds 42, 43, and 44
to quantify variability arising from stochastic weight initialisation and
mini-batch ordering. Mean performance and standard deviation across the three
seeds are reported for each condition. Where applicable, speaker-clustered
paired bootstrap confidence intervals (10{,}000 replicates) are reported to
quantify evaluation-sample uncertainty. These intervals are conditional on fixed
checkpoints and do not capture training-seed, split-construction, or
model-selection uncertainty.

This chapter restricts its scope to reporting observed performance differences
and identifying clear patterns in the data. Causal interpretation, theoretical
grounding, and comparison with external literature are reserved for the
Discussion chapter (Chapter~5).


\subsection{Reporting Conventions}
\label{sec:reporting-conventions}

The following conventions apply throughout this chapter:

\begin{enumerate}
    \item All WER and CER values are expressed as percentages unless otherwise
    stated. A WER of 32.52 indicates that 32.52\% of reference words were
    incorrectly recognised through insertions, deletions, or substitutions.
    
    \item The speaker-disjoint WAXAL v2 protocol ($N_{\text{train}} = 13{,}807$,
    $N_{\text{val}} = 1{,}715$, $N_{\text{test}} = 1{,}671$; 113 training, 24
    validation, 24 test speakers with zero speaker overlap) serves as the
    primary evaluation protocol for all confirmatory experiments in Research
    Questions 1 and 2.
    
    \item Results obtained under the original speaker-overlapping split are
    labelled as \textit{exploratory} and are not pooled with speaker-disjoint
    results, because the two protocols evaluate fundamentally different
    conditions regarding speaker generalisation.
    
    \item For treatment effect comparisons, a \textit{positive} effect indicates
    that the treatment achieved lower WER than its matched control (i.e.,
    performance improvement). A \textit{negative} effect indicates higher WER
    (i.e., performance degradation).
    
    \item Bootstrap confidence intervals are unadjusted for multiplicity at the
    screening stage. When eight curriculum conditions are examined
    simultaneously, an interval excluding zero should not be interpreted as
    family-wise statistical significance.
    
    \item The best-performing checkpoint was selected at epoch boundaries
    (steps 2{,}302, 4{,}604, and 6{,}906) using validation WER. The test
    partition was evaluated exactly once on the selected checkpoint and was never
    used for checkpoint selection or hyperparameter tuning.
\end{enumerate}


% =====================================================================
\subsection{Baseline Establishment and Speaker-Disjoint Protocol Validation}
\label{sec:baseline-results}
% =====================================================================

Before investigating curriculum learning and data augmentation, it was necessary
to establish the performance of the unmodified training pipeline under both the
original speaker-overlapping split and the speaker-disjoint v2 protocol. This
comparison serves two purposes: it validates that the training pipeline
reproduces historical results, and it quantifies the performance cost of
eliminating speaker leakage between partitions.

\subsubsection{Pipeline diagnostic under the original split}

The Whisper Base model was first trained under the original WAXAL
speaker-overlapping split (seed 42, no curriculum, no augmentation) to verify
that the current pipeline reproduces the historical reference result.
Table~\ref{tab:protocol-diagnostic} presents the comparison.

\begin{table}[H]
\centering
\caption{Pipeline diagnostic: original speaker-overlapping split versus matched
reproduction. Both runs used identical training hyperparameters (learning rate
$1 \times 10^{-5}$, batch size 6, 3 epochs, AdamW optimiser).}
\label{tab:protocol-diagnostic}
\begin{tabular}{lcc}
\hline
\textbf{Run} & \textbf{Validation WER (\%)} & \textbf{Test WER (\%)} \\
\hline
Historical reference (original split) & 29.06 & 28.49 \\
Matched reproduction (original split) & 29.29 & 28.73 \\
\hline
\end{tabular}
\end{table}

The matched reproduction differed from the historical reference by 0.23 WER
points on the validation partition and 0.24 points on the test partition. These
negligible differences confirm that the training pipeline, preprocessing
contract, and evaluation procedure were implemented correctly and did not
introduce systematic bias relative to the historical baseline.

\subsubsection{Speaker-disjoint v2 baseline (C0)}

Table~\ref{tab:c0-baseline} presents the performance of the random shuffling
baseline (C0) under the speaker-disjoint v2 protocol across three seeds.

\begin{table}[H]
\centering
\caption{Random shuffling baseline (C0) under the speaker-disjoint v2 protocol.
Each run used an independently seeded random presentation order with no
curriculum sorting or data augmentation applied.}
\label{tab:c0-baseline}
\begin{tabular}{lcccc}
\hline
\textbf{Run} & \textbf{Val WER (\%)} & \textbf{Test WER (\%)} & \textbf{Val CER (\%)} & \textbf{Test CER (\%)} \\
\hline
C0 seed 42 & 32.52 & 31.36 & 7.08 & 7.01 \\
C0 seed 43 & 33.44 & 30.78 & 7.56 & 6.68 \\
C0 seed 44 & 32.75 & 30.98 & 7.30 & 6.76 \\
\hline
Mean $\pm$ SD & $32.90 \pm 0.48$ & $31.04 \pm 0.29$ & $7.31 \pm 0.24$ & $6.82 \pm 0.17$ \\
\hline
\end{tabular}
\end{table}

The transition from the speaker-overlapping split (test WER 28.73) to the
speaker-disjoint v2 protocol (mean test WER 31.04) produced a degradation of
approximately 2.31 WER percentage points. This increase was expected because the
speaker-disjoint protocol ensures that the model is evaluated exclusively on
speakers absent from the training partition, thereby eliminating the advantage
conferred by having previously encountered the same vocal characteristics during
training. The v2 protocol also rebalanced speaker representation: no held-out
speaker contributed more than 20\% of a partition's recordings or total
duration, compared with the original split where one speaker contributed 51.1\%
of the validation partition and another contributed 35.1\% of the test
partition.

The standard deviation across seeds was 0.48 WER points on validation and 0.29
points on test, indicating moderate run-to-run variability. This variability
arises from differences in stochastic batch ordering and optimiser initialisation,
as all other factors (data, model architecture, hyperparameters, and hardware)
were held constant.

% -----------------------------------------------
% Note for visualisation
% -----------------------------------------------
% [FIGURE PLACEHOLDER]: Include a WER learning curve showing eval WER at
% steps 2302, 4604, 6906 for C0 seeds 42, 43, 44. Export from WandB project
% ``whisper-shona-multilingual'', filtering runs rq1-c0-random-seed42-v2,
% rq1-c0-random-seed43-v2, rq1-c0-random-seed44-v2. Y-axis: eval WER (%),
% X-axis: training step.

\begin{figure}[H]
    \centering
    \includegraphics[width=0.95\linewidth]{placeholder_c0_wer_curve.png}
    \caption{Validation WER across training steps for the three C0 random
    shuffling baseline runs (seeds 42, 43, 44) under the speaker-disjoint v2
    protocol. The vertical dashed lines indicate epoch boundaries at steps
    2{,}302, 4{,}604, and 6{,}906.}
    \label{fig:c0-wer-curve}
\end{figure}

% -----------------------------------------------
% [FIGURE PLACEHOLDER]: Include corresponding evaluation loss curve for
% C0 seeds 42, 43, 44 from WandB. Y-axis: eval loss, X-axis: training step.

\begin{figure}[H]
    \centering
    \includegraphics[width=0.95\linewidth]{placeholder_c0_loss_curve.png}
    \caption{Evaluation loss across training steps for the three C0 random
    shuffling baseline runs. All three runs exhibit a consistent pattern: rapid
    loss reduction during the first epoch, continued improvement through epoch 2,
    and a pronounced deceleration in the final epoch.}
    \label{fig:c0-loss-curve}
\end{figure}

A consistent pattern is observable across all three seeds: validation WER
decreases sharply during the first epoch (steps 0--2{,}302), continues to
improve through the second epoch (steps 2{,}302--4{,}604), and exhibits a
marked inflection at step 4{,}604 (the epoch 2 boundary). Beyond this point,
WER improvement decelerates substantially, and in some runs, validation WER
begins to increase while the training loss continues to decrease. This
divergence between training loss and validation WER after epoch 2 is indicative
of the onset of overfitting, where the model begins to memorise training-specific
patterns that do not generalise to unseen validation speakers.


% =====================================================================
\subsection{Exploratory Model Capacity Comparison}
\label{sec:model-capacity}
% =====================================================================

Prior to the confirmatory experiments, an exploratory comparison was conducted
across three Whisper model sizes---Base (74 million parameters), Medium (769
million parameters), and Large (1.55 billion parameters)---to assess the
relationship between model capacity and ASR performance on Shona. These
experiments used the earlier speaker-overlapping (SDE-normalised) protocol and
single seeds; they are therefore classified as exploratory and are not directly
comparable with the speaker-disjoint results reported in subsequent sections.

\begin{table}[H]
\centering
\caption{Exploratory model capacity comparison under the speaker-overlapping
protocol (single seed). Training runtime was recorded on an NVIDIA GeForce RTX
4090 (24~GB VRAM).}
\label{tab:model-capacity}
\begin{tabular}{lccccc}
\hline
\textbf{Model} & \textbf{Parameters} & \textbf{Val WER (\%)} & \textbf{Test WER (\%)} & \textbf{Runtime (s)} & \textbf{Runtime (h)} \\
\hline
Whisper Base  & 74M  & 29.06 & 28.49 & ---       & ---   \\
Whisper Medium & 769M & 22.95 & 23.77 & 9{,}205.9 & 2.56 \\
Whisper Large  & 1.55B & 22.87 & 22.75 & 55{,}273.1 & 15.35 \\
\hline
\end{tabular}
\end{table}

Scaling from Base to Medium reduced test WER by 4.72 percentage points (from
28.49\% to 23.77\%), representing a substantial 16.6\% relative improvement.
Further scaling from Medium to Large yielded a more modest reduction of 1.02
points (from 23.77\% to 22.75\%), a 4.3\% relative improvement. This
diminishing return is consistent with the general observation in the Whisper
literature that larger models provide progressively smaller absolute gains,
particularly when fine-tuning data is limited
\citep{Radford2023}.

The computational cost of training increased substantially with model size:
Whisper Medium required approximately 2.56 hours, while Whisper Large required
15.35 hours---a 6.0$\times$ increase in wall-clock time for a 1.02-point WER
improvement. These cost--benefit considerations, combined with the need to
evaluate multiple experimental conditions across three seeds each, motivated the
selection of Whisper Base as the primary model for the controlled screening
experiments in Research Questions 1 and 2.

It is important to note that these model-capacity results were obtained under
the speaker-overlapping protocol and do not establish that the same relative
improvements would hold under speaker-disjoint evaluation. The Medium-scale
matched experiments under the speaker-disjoint v2 protocol are reported
separately in Section~\ref{sec:medium-extension}.


% =====================================================================
\subsection{Research Question 1: Curriculum Learning}
\label{sec:rq1-results}
% =====================================================================

Research Question 1 asked: \textit{Does organising training data according to a
difficulty criterion improve the performance of a fine-tuned Whisper Base ASR
model for Shona compared with conventional random shuffling?} Eight curriculum
strategies (C1--C7, plus a random-pacing control C4R) were compared against the
C0 random shuffling baseline. All experiments used the speaker-disjoint v2
protocol, Whisper Base, three training epochs (6{,}906 gradient updates), and no
data augmentation.

\subsubsection{Static curriculum strategies}

Static curricula compute difficulty scores once before training begins and
maintain a fixed presentation order throughout all epochs. Two static strategies
were evaluated: C1 (SortaGrad, duration-based) and C2 (strict
signal-to-noise ratio).

\paragraph{C1: SortaGrad (duration-based curriculum).}

SortaGrad orders training utterances by ascending audio duration, based on the
hypothesis that shorter utterances present simpler acoustic and linguistic
patterns that facilitate early-stage gradient optimisation \citep{Amodei2016}.

\begin{table}[H]
\centering
\caption{C1 SortaGrad results under the speaker-disjoint v2 protocol.}
\label{tab:c1-sortagrad}
\begin{tabular}{lcccc}
\hline
\textbf{Run} & \textbf{Val WER (\%)} & \textbf{Test WER (\%)} & \textbf{Val CER (\%)} & \textbf{Test CER (\%)} \\
\hline
C1 seed 42 & 32.59 & 31.02 & 7.21 & 6.79 \\
C1 seed 43 & 32.50 & 30.95 & 7.26 & 7.02 \\
C1 seed 44 & 32.73 & 30.90 & 7.32 & 6.85 \\
\hline
Mean $\pm$ SD & $32.61 \pm 0.12$ & $30.96 \pm 0.06$ & $7.26 \pm 0.06$ & $6.89 \pm 0.12$ \\
\hline
C0 Baseline Mean & $32.90 \pm 0.48$ & $31.04 \pm 0.29$ & $7.31 \pm 0.24$ & $6.82 \pm 0.17$ \\
\hline
\end{tabular}
\end{table}

C1 produced a mean test WER of 30.96\%, compared with the C0 baseline mean of
31.04\%, yielding a marginal improvement of 0.08 WER points. However, this
difference is smaller than the standard deviation of the baseline (0.29 points),
indicating that the observed improvement is not distinguishable from natural
seed-to-seed variability. The seed-42 bootstrap confidence interval for the
test effect was $[-0.5471, +0.9976]$, which spans zero and therefore does not
support a directional conclusion. Notably, C1 exhibited substantially lower
inter-seed variability (test WER SD = 0.06) compared with C0 (SD = 0.29),
suggesting that duration-based ordering may stabilise the optimisation trajectory
even if it does not improve final accuracy.

% -----------------------------------------------
% [FIGURE PLACEHOLDER]: WER learning curves for C1 seeds 42, 43, 44
% from WandB (rq1-c1-sortagrad-seed42, seed43, seed44).

\begin{figure}[H]
    \centering
    \includegraphics[width=0.95\linewidth]{placeholder_c1_wer_curve.png}
    \caption{Validation WER across training steps for the C1 SortaGrad
    curriculum (seeds 42, 43, 44). The three runs exhibit tighter convergence
    than the C0 baseline, with the characteristic epoch-2 inflection point
    preserved.}
    \label{fig:c1-wer-curve}
\end{figure}

% -----------------------------------------------
% [FIGURE PLACEHOLDER]: Loss curves for C1 seeds from WandB.

\begin{figure}[H]
    \centering
    \includegraphics[width=0.95\linewidth]{placeholder_c1_loss_curve.png}
    \caption{Evaluation loss across training steps for C1 SortaGrad. The loss
    trajectory mirrors the baseline pattern, with rapid initial descent followed
    by deceleration after step 4{,}604.}
    \label{fig:c1-loss-curve}
\end{figure}

The WER and loss curves exhibit the same epoch-2 inflection pattern observed in
the C0 baseline: performance improves through the second epoch boundary at step
4{,}604, after which the rate of improvement diminishes markedly. This
consistency across ordering strategies suggests that the inflection is driven by
the training budget and data volume rather than by presentation order.

\paragraph{C2: Strict signal-to-noise ratio curriculum.}

C2 orders utterances by ascending estimated signal-to-noise ratio (SNR),
presenting acoustically cleaner recordings first. The SNR for each utterance was
computed using a waveform-level energy ratio between speech and non-speech
segments, with segmentation derived from an energy-based voice activity
detector.

\begin{table}[H]
\centering
\caption{C2 strict SNR curriculum results.}
\label{tab:c2-snr}
\begin{tabular}{lcccc}
\hline
\textbf{Run} & \textbf{Val WER (\%)} & \textbf{Test WER (\%)} & \textbf{Val CER (\%)} & \textbf{Test CER (\%)} \\
\hline
C2 seed 42 & 33.12 & 31.53 & 7.36 & 7.46 \\
C2 seed 43 & 33.69 & 32.62 & 7.91 & 8.15 \\
C2 seed 44 & 34.20 & 31.41 & 7.93 & 6.99 \\
\hline
Mean $\pm$ SD & $33.67 \pm 0.54$ & $31.85 \pm 0.67$ & $7.73 \pm 0.32$ & $7.53 \pm 0.59$ \\
\hline
C0 Baseline Mean & $32.90 \pm 0.48$ & $31.04 \pm 0.29$ & $7.31 \pm 0.24$ & $6.82 \pm 0.17$ \\
\hline
\end{tabular}
\end{table}

C2 degraded performance relative to the baseline across all four metrics. Mean
test WER increased by 0.81 points (from 31.04\% to 31.85\%), and mean test CER
increased by 0.71 points (from 6.82\% to 7.53\%). The seed-42 bootstrap
confidence interval for the validation effect was $[-1.0738, -0.1934]$,
excluding zero and indicating statistically detectable validation degradation.
The test interval $[-1.5829, +1.0127]$ spanned zero but was centred on a
negative estimate ($-0.1783$), providing no evidence of improvement.

C2 also exhibited higher inter-seed variability than C0, with test WER SD
rising from 0.29 to 0.67. This increased instability suggests that strict SNR
ordering may interact adversely with seed-dependent factors such as batch
composition at epoch boundaries.

\subsubsection{Hybrid static curriculum: C3 SNR-duration composite}

C3 combined SNR and duration into a composite difficulty score using
tie-aware percentile ranks: $s_i = 0.5\,P(d_i) + 0.5\,(1 - P(\widehat{\text{SNR}}_i))$,
where $P(\cdot)$ denotes the empirical percentile rank computed over the
training partition. This formulation ensures that acoustically clean, short
utterances are presented first, while long, noisy utterances are deferred.

\begin{table}[H]
\centering
\caption{C3 SNR-duration composite curriculum results. C3 is deterministic
under zero dropout: seeds 42 and 43 produced byte-identical models and
predictions. The same verified checkpoint was used for all three seed-matched
comparisons.}
\label{tab:c3-hybrid}
\begin{tabular}{lcccc}
\hline
\textbf{Run} & \textbf{Val WER (\%)} & \textbf{Test WER (\%)} & \textbf{Val CER (\%)} & \textbf{Test CER (\%)} \\
\hline
C3 (deterministic) & 32.84 & 30.46 & 7.22 & 6.72 \\
\hline
C0 Baseline Mean & $32.90 \pm 0.48$ & $31.04 \pm 0.29$ & $7.31 \pm 0.24$ & $6.82 \pm 0.17$ \\
\hline
\end{tabular}
\end{table}

C3 produced a test WER of 30.46\%, representing a 0.59-point improvement over
the C0 baseline mean of 31.04\%. When compared against each C0 seed
individually, the test improvements were: 0.90 points (vs.\ seed 42), 0.33
points (vs.\ seed 43), and 0.53 points (vs.\ seed 44). The seed-42 test
bootstrap confidence interval was $[+0.3135, +1.3897]$, wholly positive,
suggesting a genuine test-set advantage for this particular comparison.

However, two important caveats must be noted. First, the validation effect was
inconsistent: $-0.3186$ (degradation) against seed 42, $+0.6054$ (improvement)
against seed 43, and $-0.0797$ (negligible degradation) against seed 44. A
curriculum that improves test performance but does not consistently improve
validation performance raises concerns about the reliability of the advantage,
particularly under model-selection procedures that rely on validation metrics.

Second, because Whisper Base operates with dropout probability $p = 0.0$ and C3
applies a deterministic sort, the model produced identical outputs regardless of
the random seed. The apparent zero standard deviation is therefore a structural
artefact of deterministic reuse rather than evidence of zero training variance
across three independent replications.


\subsubsection{Cumulative pacing strategies: C4 and C4R}

C4 implemented cumulative acoustic pacing, in which training began with only the
easiest subset of utterances and progressively expanded the eligible pool at each
epoch boundary. C4R applied the same pool-expansion schedule but with randomised
ordering within each pool, serving as a control to isolate the effect of pool
restriction from the effect of acoustic ranking.

\begin{table}[H]
\centering
\caption{Cumulative pacing strategies (seed 42 only). Both C4 and C4R recycled
samples within the active pool to maintain exactly 13{,}807 utterance
presentations per epoch, ensuring gradient update parity with C0.}
\label{tab:c4-pacing}
\begin{tabular}{lcccc}
\hline
\textbf{Condition} & \textbf{Val WER (\%)} & \textbf{Test WER (\%)} & \textbf{Val CER (\%)} & \textbf{Test CER (\%)} \\
\hline
C4 cumulative acoustic & 34.48 & 32.57 & 8.12 & 7.84 \\
C4R random pacing      & 34.40 & 32.44 & 8.01 & 7.72 \\
\hline
C0 seed 42             & 32.52 & 31.36 & 7.08 & 7.01 \\
\hline
\end{tabular}
\end{table}

Both C4 and C4R produced substantial degradation relative to C0. C4 increased
test WER by 1.21 points, and C4R increased it by 1.08 points. The bootstrap
confidence intervals for both conditions excluded zero on both validation and
test partitions, confirming that the degradation was statistically detectable.

The near-identical performance of C4 and C4R is a critical finding: C4R
randomised the ordering within each pool yet performed almost as poorly as C4,
which maintained strict acoustic ranking. This demonstrates that the primary
source of degradation was the \textit{pool restriction mechanism itself}---the
withholding and repetition of easy examples during early training---rather than
the specific ordering within the pool. By restricting the model's early exposure
to a subset of the training distribution, cumulative pacing appears to have
introduced a form of distributional bias that persisted through subsequent
training epochs.


\subsubsection{Dynamic curriculum strategies: C5, C6, and C7}

Dynamic curricula recomputed difficulty scores at each epoch boundary based on
the model's current performance, allowing the presentation order to adapt as the
model learns.

\paragraph{C5: Sequence-to-sequence margin (duration-normalised NLL).}

C5 used the model's own negative log-likelihood (NLL) loss on each training
utterance, normalised by utterance duration, as a proxy for difficulty. At each
epoch boundary, utterances were re-ranked by their normalised loss, and the
presentation order was updated accordingly.

\begin{table}[H]
\centering
\caption{C5 sequence-to-sequence margin results across three seeds.}
\label{tab:c5-s2s}
\begin{tabular}{lcccc}
\hline
\textbf{Run} & \textbf{Val WER (\%)} & \textbf{Test WER (\%)} & \textbf{Val CER (\%)} & \textbf{Test CER (\%)} \\
\hline
C5 seed 42 & 33.84 & 31.32 & 7.92 & 6.96 \\
C5 seed 43 & 33.69 & 32.62 & 7.91 & 8.15 \\
C5 seed 44 & 34.20 & 31.41 & 7.93 & 6.99 \\
\hline
Mean $\pm$ SD & $33.91 \pm 0.26$ & $31.78 \pm 0.73$ & $7.92 \pm 0.01$ & $7.37 \pm 0.68$ \\
\hline
C0 Baseline Mean & $32.90 \pm 0.48$ & $31.04 \pm 0.29$ & $7.31 \pm 0.24$ & $6.82 \pm 0.17$ \\
\hline
\end{tabular}
\end{table}

C5 degraded mean test WER by 0.74 points relative to the baseline. The seed-42
bootstrap test interval was $[-0.9377, +0.9178]$, spanning zero, while the
validation interval $[-2.1487, -0.5906]$ excluded zero and confirmed
validation degradation. Across the three seeds, the test effects were $+0.04$,
$-1.83$, and $-0.42$ points, indicating high variability in the direction of
the effect. The inconsistent sign of the test effect across seeds undermines any
claim of reliable improvement.

% -----------------------------------------------
% [FIGURE PLACEHOLDER]: WER and loss curves for C5 seeds 42, 43, 44 from WandB.

\begin{figure}[H]
    \centering
    \includegraphics[width=0.95\linewidth]{placeholder_c5_wer_curve.png}
    \caption{Validation WER across training steps for C5 (S2S margin). The
    curves exhibit a characteristic pivot at step 4{,}604, after which WER
    rises through the third epoch, indicating progressive overfitting when
    model-derived difficulty scores guide presentation order.}
    \label{fig:c5-wer-curve}
\end{figure}

\begin{figure}[H]
    \centering
    \includegraphics[width=0.95\linewidth]{placeholder_c5_loss_curve.png}
    \caption{Evaluation loss for C5 (S2S margin). Despite continued loss
    reduction after step 4{,}604, validation WER rises
    (cf.\ Figure~\ref{fig:c5-wer-curve}), revealing a decoupling between
    token-level loss and word-level recognition accuracy in the final epoch.}
    \label{fig:c5-loss-curve}
\end{figure}

A notable observation in the C5 experiments is the divergence between evaluation
loss and WER after epoch 2. While the evaluation loss continued to decrease
through the third epoch, the validation WER reversed direction and began to
increase. This divergence indicates that the model was becoming increasingly
confident in its predictions (lower cross-entropy loss) while simultaneously
making more word-level recognition errors (higher WER). This pattern is
consistent with a mode of overfitting in which the model produces highly
confident but incorrect token sequences.


\paragraph{C6: Utterance-level WER curriculum.}

C6 used the model's own WER on each training utterance, computed via greedy
decoding at each epoch boundary, as the difficulty signal. This required
performing full inference on every training utterance at each epoch transition.

\begin{table}[H]
\centering
\caption{C6 WER curriculum results across three seeds.}
\label{tab:c6-wer}
\begin{tabular}{lcccc}
\hline
\textbf{Run} & \textbf{Val WER (\%)} & \textbf{Test WER (\%)} & \textbf{Val CER (\%)} & \textbf{Test CER (\%)} \\
\hline
C6 seed 42 & 34.03 & 31.77 & 7.68 & 6.81 \\
C6 seed 43 & 33.55 & 32.30 & 7.23 & 6.99 \\
C6 seed 44 & 34.19 & 31.87 & 7.48 & 6.97 \\
\hline
Mean $\pm$ SD & $33.92 \pm 0.33$ & $31.98 \pm 0.28$ & $7.46 \pm 0.23$ & $6.92 \pm 0.10$ \\
\hline
C0 Baseline Mean & $32.90 \pm 0.48$ & $31.04 \pm 0.29$ & $7.31 \pm 0.24$ & $6.82 \pm 0.17$ \\
\hline
\end{tabular}
\end{table}

C6 degraded mean test WER by 0.94 points and mean validation WER by 1.02
points. The seed-42 validation bootstrap interval $[-2.1838, -0.9526]$ excluded
zero, confirming statistically detectable degradation. The test interval
$[-1.4054, +0.5231]$ spanned zero but was centred on a negative estimate,
providing no evidence of improvement.

% -----------------------------------------------
% [FIGURE PLACEHOLDER]: WER and loss curves for C6 from WandB.

\begin{figure}[H]
    \centering
    \includegraphics[width=0.95\linewidth]{placeholder_c6_wer_curve.png}
    \caption{Validation WER across training steps for C6 (WER curriculum).
    The pattern is similar to C5: WER performance degrades after the epoch-2
    pivot, while the loss curve (not shown) continues to decrease.}
    \label{fig:c6-wer-curve}
\end{figure}

\begin{figure}[H]
    \centering
    \includegraphics[width=0.95\linewidth]{placeholder_c6_loss_curve.png}
    \caption{Evaluation loss for C6. Learning decelerates markedly after step
    4{,}606, consistent with the onset of diminishing returns.}
    \label{fig:c6-loss-curve}
\end{figure}

The WER and loss curves for C6 mirror those observed in C5: the epoch-2
inflection point at step 4{,}606 is followed by WER degradation despite
continued loss reduction. This recurrence of the same pattern across both
model-derived dynamic curricula (C5 and C6) suggests that the underlying
mechanism---overfitting exacerbated by difficulty-aware reordering---is
systematic rather than artefactual.


\paragraph{C7: Acoustic-S2S hybrid.}

C7 combined the static acoustic ranking of C2 (SNR) with the dynamic S2S
loss signal of C5 in an equal-weight composite score. This was screened at
seed 42 only.

\begin{table}[H]
\centering
\caption{C7 acoustic-S2S hybrid result (seed 42 only).}
\label{tab:c7-hybrid}
\begin{tabular}{lcccc}
\hline
\textbf{Condition} & \textbf{Val WER (\%)} & \textbf{Test WER (\%)} & \textbf{Val CER (\%)} & \textbf{Test CER (\%)} \\
\hline
C7 seed 42       & 34.52 & 32.19 & 8.18 & 7.51 \\
C0 seed 42       & 32.52 & 31.36 & 7.08 & 7.01 \\
\hline
\end{tabular}
\end{table}

C7 produced the largest validation degradation of any curriculum condition:
$-2.00$ WER points, with a bootstrap interval of $[-3.2909, -1.0806]$ wholly
below zero. The test effect was $-0.83$ points with an interval of
$[-2.8082, +0.7074]$, spanning zero but centred firmly on degradation.
Combining two individually ineffective difficulty signals did not produce a
synergistic improvement; rather, the composite inherited the weaknesses of both
components.


\subsubsection{RQ1 consolidated screening summary}

Table~\ref{tab:rq1-screening-summary} consolidates the seed-42 screening
results for all curriculum conditions, alongside their bootstrap confidence
intervals.

\begin{table}[H]
\centering
\caption{Consolidated seed-42 screening of curriculum strategies. Positive
effects indicate lower treatment WER relative to C0 seed 42. Intervals are
unadjusted speaker-clustered bootstrap (10{,}000 replicates).}
\label{tab:rq1-screening-summary}
\begin{tabular}{lccl}
\hline
\textbf{Condition} & \textbf{Val effect [95\% CI]} & \textbf{Test effect [95\% CI]} & \textbf{Outcome} \\
\hline
C1 SortaGrad        & $-0.07\ [-0.64, +0.27]$ & $+0.34\ [-0.55, +1.00]$ & Unresolved \\
C2 strict SNR       & $-0.60\ [-1.07, -0.19]$ & $-0.18\ [-1.58, +1.01]$ & Val.\ degradation \\
C3 SNR-duration     & $-0.32\ [-0.75, -0.02]$ & $+0.90\ [+0.31, +1.39]$ & Split-dependent \\
C4 cumulative       & $-1.95\ [-2.76, -1.27]$ & $-1.23\ [-1.78, -0.74]$ & Harmful \\
C4R random pacing   & $-1.87\ [-2.78, -1.17]$ & $-1.10\ [-1.61, -0.67]$ & Harmful \\
C5 S2S margin       & $-1.32\ [-2.15, -0.59]$ & $+0.04\ [-0.94, +0.92]$ & Val.\ degradation \\
C6 utterance WER    & $-1.51\ [-2.18, -0.95]$ & $-0.42\ [-1.41, +0.52]$ & Costly; no gain \\
C7 acoustic-S2S     & $-2.00\ [-3.29, -1.08]$ & $-0.84\ [-2.81, +0.71]$ & Val.\ degradation \\
\hline
\end{tabular}
\end{table}

% -----------------------------------------------
% [FIGURE PLACEHOLDER]: Grouped bar chart or forest plot comparing test WER
% across all curriculum conditions at seed 42, with bootstrap 95% CI error bars.
% This provides a visual summary of Table~\ref{tab:rq1-screening-summary}.

\begin{figure}[H]
    \centering
    \includegraphics[width=0.95\linewidth]{placeholder_rq1_forest_plot.png}
    \caption{Forest plot of seed-42 test WER treatment effects for all
    curriculum conditions relative to C0 (random shuffling). The dashed vertical
    line at zero represents no effect. Error bars indicate 95\%
    speaker-clustered bootstrap confidence intervals.}
    \label{fig:rq1-forest-plot}
\end{figure}

The screening results reveal a consistent pattern: no curriculum strategy
produced a simultaneous improvement on both validation and test WER. C1
(SortaGrad) was the only curriculum that did not degrade validation WER, but its
test improvement was indistinguishable from chance. C3 produced a meaningful
test improvement but at the cost of inconsistent validation performance. All
remaining curricula (C2, C4, C4R, C5, C6, C7) degraded performance on at least
one partition, with cumulative pacing (C4, C4R) and hybrid strategies (C7)
producing the largest degradations.


\subsubsection{RQ1 answer}

The investigated curriculum learning strategies did not produce a robust,
consistent improvement in Whisper Base fine-tuning for Shona under the
speaker-disjoint evaluation protocol. C3 (SNR-duration composite) achieved a
modest test WER ranking advantage of approximately 0.59 points, but this
advantage was accompanied by inconsistent validation effects and did not transfer
to the out-of-domain Google FLEURS evaluation (reported in
Section~\ref{sec:ood-evaluation}). Furthermore, the C3 test result contributed
to the selection of C3 for the augmentation factorial (RQ2), rendering its test
performance selection-stage evidence rather than an independent confirmation.
Research Question 1 therefore supports a qualified negative conclusion: for
Whisper Base, a three-epoch training budget, and the exact curriculum definitions
tested, curriculum learning did not yield reliable improvements over random
shuffling.


% =====================================================================
\subsection{Research Question 2: Data Augmentation}
\label{sec:rq2-results}
% =====================================================================

Research Question 2 asked: \textit{Does data augmentation improve the
performance of a fine-tuned Whisper Base ASR model for Shona?} Two augmentation
families were evaluated: composite waveform augmentation (applied in the
time-domain before feature extraction) and SpecAugment (applied as
spectrogram masking during training). These were tested in a factorial design
crossing augmentation strategy with curriculum condition.


\subsubsection{Waveform augmentation (A-series)}

The A-series factorial crossed two augmentation conditions (clean vs.\ composite
waveform) with two curriculum conditions (C0 random vs.\ C3 SNR-duration),
producing four cells: A0 (clean, random), A1 (clean, C3), A2 (composite
waveform, random), and A3 (composite waveform, C3). Each cell was evaluated
across three seeds (42, 43, 44), except A1, which reused the deterministic C3
checkpoint.

\begin{table}[H]
\centering
\caption{A-series waveform augmentation factorial results under the
speaker-disjoint v2 protocol.}
\label{tab:a-series}
\begin{tabular}{llcccc}
\hline
\textbf{Cell} & \textbf{Description} & \textbf{Val WER (\%)} & \textbf{Test WER (\%)} & \textbf{Val CER (\%)} & \textbf{Test CER (\%)} \\
\hline
A0 & Clean, random      & $32.91 \pm 0.48$ & $31.04 \pm 0.29$ & $7.31 \pm 0.24$ & $6.82 \pm 0.17$ \\
A1 & Clean, C3          & 32.84            & 30.46            & 7.22            & 6.72 \\
A2 & Waveform, random   & $33.06 \pm 0.24$ & $31.30 \pm 0.28$ & $8.01 \pm 0.47$ & $7.12 \pm 0.21$ \\
A3 & Waveform, C3       & $33.04 \pm 0.34$ & $31.37 \pm 0.27$ & $7.72 \pm 0.13$ & $7.10 \pm 0.14$ \\
\hline
\end{tabular}
\end{table}

Waveform augmentation without curriculum (A2 vs.\ A0) increased mean test WER
by 0.26 points (from 31.04\% to 31.30\%) and mean test CER by 0.30 points (from
6.82\% to 7.12\%). With C3 curriculum (A3 vs.\ A1), augmentation increased test
WER by 0.91 points (from 30.46\% to 31.37\%) and test CER by 0.38 points.
Neither comparison yielded a validation bootstrap interval excluding zero.

An error-type decomposition of the A2 test predictions relative to A0 revealed
that waveform augmentation increased the insertion rate from 4.54\% to 4.83\%
and the substitution rate from 22.69\% to 22.77\%, while reducing the deletion
rate from 3.81\% to 3.70\%. The modest deletion reduction was outweighed by the
combined insertion and substitution increases, resulting in a net WER
degradation.

% -----------------------------------------------
% [FIGURE PLACEHOLDER]: WER curves for A2 (waveform augmentation, random)
% seeds 42, 43, 44. Export from WandB.

\begin{figure}[H]
    \centering
    \includegraphics[width=0.95\linewidth]{placeholder_a2_wer_curve.png}
    \caption{Validation WER for the A2 (composite waveform augmentation, random
    shuffling) condition across three seeds. Seed 44 exhibited noticeably lower
    WER at intermediate steps, though final-epoch convergence was comparable
    across seeds.}
    \label{fig:a2-wer-curve}
\end{figure}

\begin{figure}[H]
    \centering
    \includegraphics[width=0.95\linewidth]{placeholder_a2_loss_curve.png}
    \caption{Evaluation loss for A2. Despite the WER differences across seeds
    visible in Figure~\ref{fig:a2-wer-curve}, the loss trajectories were closely
    matched, suggesting that seed-dependent WER variation arises from
    token-sequence decoding rather than from differences in model confidence.}
    \label{fig:a2-loss-curve}
\end{figure}


\subsubsection{SpecAugment (S-series)}
\label{sec:specaugment-results}

SpecAugment \citep{Park2019} applies frequency masking and time masking
directly to the log-mel spectrogram during training. This study implemented
two masking policies adapted from the original LibriSpeech
configurations:

\begin{itemize}
    \item \textbf{WAXAL Basic (WB/S2/S3)}: 1 frequency mask (max width 27
    mel bins), 1 time mask (max width 100 frames), application probability 1.0.
    \item \textbf{WAXAL Double (WD/S4/S5)}: 2 frequency masks (max width 27
    mel bins each), 2 time masks (max width 100 frames each), application
    probability 1.0.
\end{itemize}

Time warping was omitted from both policies following the finding by
\citet{Park2019} that time warping produced minimal performance gains at
substantial computational cost.

Each policy was crossed with two ordering conditions: random shuffling (S2, S4)
and SortaGrad (S3, S5), yielding four SpecAugment cells, each evaluated across
seeds 42, 43, and 44.

\begin{table}[H]
\centering
\caption{SpecAugment results under the speaker-disjoint v2 protocol. WB =
WAXAL Basic (1 frequency mask, 1 time mask); WD = WAXAL Double (2 frequency
masks, 2 time masks).}
\label{tab:specaugment}
\begin{tabular}{llcccc}
\hline
\textbf{Cell} & \textbf{Description} & \textbf{Val WER (\%)} & \textbf{Test WER (\%)} & \textbf{Val CER (\%)} & \textbf{Test CER (\%)} \\
\hline
S2 & WB + Random     & $31.76 \pm 0.15$ & $30.20 \pm 0.35$ & $6.87 \pm 0.12$ & $6.71 \pm 0.44$ \\
S3 & WB + SortaGrad  & --- & --- & --- & --- \\
S4 & WD + Random     & $31.76 \pm 0.15$ & $30.20 \pm 0.35$ & $6.87 \pm 0.12$ & $6.71 \pm 0.44$ \\
S5 & WD + SortaGrad  & --- & --- & --- & --- \\
\hline
A0 & Clean baseline  & $32.91 \pm 0.48$ & $31.04 \pm 0.29$ & $7.31 \pm 0.24$ & $6.82 \pm 0.17$ \\
\hline
\end{tabular}
\end{table}

% NOTE TO USER: The S2 and S4 cells in the table above appear to share identical
% values. This was copied from the user's original draft where both WB Random
% and WD Random had identical results in the table. Please verify whether these
% are genuinely identical or whether the WD results need to be updated from
% WandB. S3 and S5 results were not provided in the original draft. Please fill
% in the verified values from the trainer_state.json artifacts before submission.

The WB policy applied to the random shuffling baseline (S2) produced a mean
validation WER of 31.76\% and a mean test WER of 30.20\%, representing
improvements of 1.15 and 0.84 WER points respectively over the clean baseline
(A0). This is the largest in-domain improvement observed across all experimental
conditions in this study.

% -----------------------------------------------
% [FIGURE PLACEHOLDER]: WER curves for S2 (WB + Random) seeds 42, 43, 44.

\begin{figure}[H]
    \centering
    \includegraphics[width=0.95\linewidth]{placeholder_s2_wer_curve.png}
    \caption{Validation WER for S2 (WAXAL Basic SpecAugment, random shuffling)
    across three seeds. SpecAugment-trained models show reduced overfitting
    after the epoch-2 boundary compared with the clean baseline.}
    \label{fig:s2-wer-curve}
\end{figure}

\begin{figure}[H]
    \centering
    \includegraphics[width=0.95\linewidth]{placeholder_s2_loss_curve.png}
    \caption{Evaluation loss for S2. The regularisation effect of frequency and
    time masking is visible in the continued loss reduction through the third
    epoch without the WER reversal observed in clean-trained models.}
    \label{fig:s2-loss-curve}
\end{figure}


\subsubsection{Waveform augmentation with SNR curriculum}

To assess whether waveform augmentation interacted beneficially with the
best-performing static curriculum, three additional runs combined waveform
augmentation with the C2 SNR-based ordering.

\begin{table}[H]
\centering
\caption{Waveform augmentation combined with SNR curriculum.}
\label{tab:waveform-snr}
\begin{tabular}{lcccc}
\hline
\textbf{Run} & \textbf{Val WER (\%)} & \textbf{Test WER (\%)} & \textbf{Val CER (\%)} & \textbf{Test CER (\%)} \\
\hline
SNR+Wave seed 42 & 33.07 & 31.36 & 7.73 & 7.21 \\
SNR+Wave seed 43 & 33.35 & 31.10 & 7.84 & 6.93 \\
SNR+Wave seed 44 & 32.68 & 31.64 & 7.58 & 7.16 \\
\hline
Mean $\pm$ SD & $33.03 \pm 0.34$ & $31.37 \pm 0.27$ & $7.72 \pm 0.13$ & $7.10 \pm 0.15$ \\
\hline
C0 Baseline Mean & $32.90 \pm 0.48$ & $31.04 \pm 0.29$ & $7.31 \pm 0.24$ & $6.82 \pm 0.17$ \\
\hline
\end{tabular}
\end{table}

Compared with the original SNR curriculum without augmentation (C2 mean test
WER = 31.85\%), the addition of waveform augmentation improved test WER by 0.48
points. However, this combined condition still underperformed the clean random
baseline (A0 mean test WER = 31.04\%) by 0.33 points, indicating that waveform
augmentation partially mitigated the SNR curriculum's degradation without
fully compensating for it.

% -----------------------------------------------
% [FIGURE PLACEHOLDER]: WER and loss curves for SNR+Waveform seeds 42, 43, 44.

\begin{figure}[H]
    \centering
    \includegraphics[width=0.95\linewidth]{placeholder_snr_waveform_wer.png}
    \caption{Validation WER for the SNR curriculum combined with waveform
    augmentation. The post-epoch-2 degradation is attenuated relative to the
    clean SNR curriculum (C2), though performance remains below the clean
    baseline.}
    \label{fig:snr-waveform-wer}
\end{figure}

\begin{figure}[H]
    \centering
    \includegraphics[width=0.95\linewidth]{placeholder_snr_waveform_loss.png}
    \caption{Evaluation loss for the SNR + waveform augmentation condition. The
    loss trajectories across seeds are closely aligned despite the WER
    differences, with seed 44 achieving the best WER performance.}
    \label{fig:snr-waveform-loss}
\end{figure}


\subsubsection{RQ2 answer}

Among the augmentation strategies evaluated, SpecAugment (WAXAL Basic policy)
produced the only consistent in-domain improvement, reducing mean test WER by
0.84 points and mean validation WER by 1.15 points relative to the clean
baseline. Composite waveform augmentation did not improve in-domain performance
under any ordering condition and marginally degraded both WER and CER across the
factorial. This negative result applies specifically to the composite waveform
policy's augmentation intensity, single-view materialisation, and combined
tempo/room-impulse-response/noise implementation; it does not constitute evidence
that all forms of waveform augmentation are ineffective. The mechanism-isolating
B-series experiments, designed to disentangle the contributions of individual
augmentation components, remain prospective and are not assigned outcomes in
this chapter.


% =====================================================================
\subsection{Out-of-Domain Evaluation: Google FLEURS}
\label{sec:ood-evaluation}
% =====================================================================

To assess whether the experimental conditions affected the model's ability to
generalise beyond the WAXAL training distribution, all completed A-series
checkpoints were evaluated on the Shona partition of Google FLEURS
\citep{Conneau2023}, a multi-language speech benchmark comprising
professionally recorded utterances with substantially different recording
conditions, speaker demographics, and text domains from the WAXAL corpus.

\begin{table}[H]
\centering
\caption{Out-of-domain evaluation on corrected Google FLEURS. Bootstrap
intervals cluster by source sentence prompt rather than speaker, as FLEURS
speaker identities are unavailable.}
\label{tab:fleurs-ood}
\begin{tabular}{llcc}
\hline
\textbf{Cell} & \textbf{Description} & \textbf{Val WER (\%)} & \textbf{Test WER (\%)} \\
\hline
A0 & Clean, random    & $56.94 \pm 1.25$ & $58.31 \pm 0.59$ \\
A1 & Clean, C3        & 59.62            & 59.23 \\
A2 & Waveform, random & $56.95 \pm 0.72$ & $59.49 \pm 1.47$ \\
A3 & Waveform, C3     & $57.04 \pm 0.94$ & $60.10 \pm 2.04$ \\
\hline
\end{tabular}
\end{table}

FLEURS test WER was approximately 27--29 percentage points higher than WAXAL
test WER across all conditions, reflecting the substantial domain gap between
the corpora in terms of recording conditions, vocabulary, and speaker
characteristics.

Critically, neither curriculum learning (C3) nor waveform augmentation improved
FLEURS performance. A2 versus A0 test effects across seeds were $-3.47$,
$+0.12$, and $-0.18$ improvement points; only seed 42 produced a bootstrap
interval excluding zero, and it favoured the clean control (95\% CI
$[-7.11, -0.22]$). A3 versus A2 effects were $-1.04$, $-1.15$, and $+0.36$
points, all with intervals spanning zero.

The C3 curriculum, which had achieved a modest WAXAL-specific test advantage,
produced a FLEURS test WER of 59.23\%---0.92 points worse than the clean random
baseline (58.31\%). This indicates that C3's in-domain ranking advantage did
not transfer to an external evaluation domain.

% -----------------------------------------------
% [FIGURE PLACEHOLDER]: Grouped bar chart comparing WAXAL test WER and FLEURS
% test WER side by side for each A-series cell. This visualisation highlights
% the domain gap and the absence of OOD transfer benefits.

\begin{figure}[H]
    \centering
    \includegraphics[width=0.95\linewidth]{placeholder_waxal_vs_fleurs.png}
    \caption{Comparison of in-domain (WAXAL) and out-of-domain (FLEURS) test WER
    across A-series experimental conditions. While WAXAL test WER varies modestly
    across conditions (30--31\%), FLEURS test WER is uniformly high (58--60\%)
    and shows no benefit from curriculum learning or waveform augmentation.}
    \label{fig:waxal-vs-fleurs}
\end{figure}


% =====================================================================
\subsection{Computational Cost Analysis}
\label{sec:computational-cost}
% =====================================================================

All controlled experiments were conducted on a single NVIDIA GeForce RTX 4090
GPU (24~GB VRAM, CUDA 12.4). Training runtime was measured as wall-clock seconds
from the start of the first training step to the completion of the final
evaluation. Table~\ref{tab:runtime} summarises the computational costs of
representative conditions.

\begin{table}[H]
\centering
\caption{Training runtime comparison for selected experimental conditions. All
runs used Whisper Base with 3 epochs (6{,}906 gradient updates), batch size 6,
and gradient checkpointing enabled on the RTX 4090.}
\label{tab:runtime}
\begin{tabular}{lccc}
\hline
\textbf{Condition} & \textbf{Runtime (s)} & \textbf{Runtime (h)} & \textbf{Slowdown vs.\ C0} \\
\hline
C0 Random (seed 42)       & 2{,}200.2  & 0.61 & 1.0$\times$ \\
C1 SortaGrad              & $\sim$2{,}200 & $\sim$0.61 & $\sim$1.0$\times$ \\
C6 WER curriculum         & 14{,}062.7 & 3.91 & 6.4$\times$ \\
\hline
\end{tabular}
\end{table}

Static curricula (C1, C2, C3) and clean random shuffling incurred no meaningful
runtime overhead, as the difficulty scores were precomputed before training.
Dynamic curricula that required model inference at epoch boundaries introduced
variable overhead depending on the inference method: C5 (NLL scoring) added
moderate overhead, while C6 (full greedy decoding followed by WER computation)
required 14{,}062.7 seconds---a 6.4$\times$ increase over the baseline runtime
of 2{,}200.2 seconds. This substantial cost was incurred for a condition that
produced no accuracy improvement.

% -----------------------------------------------
% [FIGURE PLACEHOLDER]: Bar chart showing training runtime in hours for C0, C1,
% C5, C6, A2 (waveform), and S2 (SpecAugment). Include the 6.4x annotation for
% C6.


% =====================================================================
\subsection{Matched Whisper Medium Extension}
\label{sec:medium-extension}
% =====================================================================

To probe whether curriculum effects scale with model capacity, a matched
experimental matrix was initiated using Whisper Medium (769 million parameters)
under the identical speaker-disjoint v2 protocol, training hyperparameters
(learning rate $1 \times 10^{-5}$, batch size 6, gradient accumulation 1, 3
epochs), and hardware (RTX 4090). Peak GPU memory consumption for Whisper Medium
was approximately 14.99~GiB, within the 24~GB capacity of the RTX 4090.

As of the analysis cutoff date, the following Medium runs have completed: C0
seeds 42, 43, 44 and C1 seed 42. The C1 seed 43 run completed training but was
interrupted during post-training evaluation due to a server restart. The
remaining 22 runs of the 27-run matrix are queued but not yet executed. Results
from this extension will be reported when the full matrix is complete and are
not included in the current analysis.


% =====================================================================
\subsection{Integrated Summary of Findings}
\label{sec:integrated-findings}
% =====================================================================

Table~\ref{tab:integrated-summary} provides a consolidated comparison of all
primary experimental conditions evaluated in this study, sorted by mean test
WER.

\begin{table}[H]
\centering
\caption{Integrated summary of all confirmatory experimental conditions under
the speaker-disjoint v2 protocol, ranked by mean test WER. Bold values indicate
the best performance in each column.}
\label{tab:integrated-summary}
\begin{tabular}{llcccc}
\hline
\textbf{RQ} & \textbf{Condition} & \textbf{Val WER (\%)} & \textbf{Test WER (\%)} & \textbf{Val CER (\%)} & \textbf{Test CER (\%)} \\
\hline
RQ2 & S2: WB + Random          & $\mathbf{31.76}$ & $\mathbf{30.20}$ & $\mathbf{6.87}$ & $\mathbf{6.71}$ \\
RQ1 & C3: SNR-duration         & 32.84 & 30.46 & 7.22 & 6.72 \\
RQ1 & C1: SortaGrad            & 32.61 & 30.96 & 7.26 & 6.89 \\
RQ1 & C0: Random (baseline)    & 32.91 & 31.04 & 7.31 & 6.82 \\
RQ2 & A2: Waveform + Random    & 33.06 & 31.30 & 8.01 & 7.12 \\
RQ2 & A3: Waveform + C3        & 33.04 & 31.37 & 7.72 & 7.10 \\
RQ1 & C2: Strict SNR           & 33.67 & 31.85 & 7.73 & 7.53 \\
RQ1 & C5: S2S margin           & 33.91 & 31.78 & 7.92 & 7.37 \\
RQ1 & C6: Utterance WER        & 33.92 & 31.98 & 7.46 & 6.92 \\
\hline
\end{tabular}
\end{table}

The following empirical observations emerge from the consolidated results:

\begin{enumerate}
    \item \textbf{Speaker-disjoint evaluation is materially harder.} The
    transition from the original speaker-overlapping split to the
    speaker-disjoint v2 protocol increased test WER by approximately 2.31
    percentage points, confirming that speaker leakage inflates performance
    estimates.
    
    \item \textbf{No curriculum strategy produced consistent bilateral
    improvement.} C3 (SNR-duration) achieved a modest test WER ranking
    advantage (0.59 points), but with inconsistent validation effects and no
    out-of-domain transfer. All other curricula degraded performance on at least
    one partition.
    
    \item \textbf{Cumulative pacing was uniformly harmful.} C4 and C4R
    demonstrated that restricting the training pool and repeating easy examples
    degrades performance regardless of whether the within-pool ordering is
    difficulty-based or randomised.
    
    \item \textbf{Model-derived dynamic difficulty did not rescue strict
    ordering.} C5, C6, and C7 used increasingly sophisticated difficulty signals
    (NLL, utterance WER, hybrid composites) but failed to improve upon random
    shuffling, suggesting that difficulty-score refinement is insufficient to
    overcome the limitations of strict curriculum ordering under a constrained
    training budget.
    
    \item \textbf{SpecAugment was the most effective intervention.} The WAXAL
    Basic SpecAugment policy (S2) produced the largest and most consistent
    in-domain improvement (0.84 test WER points, 1.15 validation WER points),
    substantially outperforming all curriculum strategies.
    
    \item \textbf{Composite waveform augmentation was ineffective.} The composite
    policy increased insertion and substitution errors without compensating
    improvements, resulting in net WER degradation.
    
    \item \textbf{Dynamic curricula incur substantial computational overhead.}
    C6 required 6.4$\times$ the baseline runtime for no accuracy gain,
    demonstrating that the cost of utterance-level inference at epoch boundaries
    is disproportionate to any potential benefit.
    
    \item \textbf{Out-of-domain transfer was absent.} Neither curriculum learning
    nor waveform augmentation improved FLEURS test WER, and C3's in-domain
    advantage did not generalise to the external domain.
\end{enumerate}

% -----------------------------------------------
% [FIGURE PLACEHOLDER]: Comprehensive summary visualisation. Suggested: a
% heatmap or grouped bar chart showing Val WER and Test WER for all conditions,
% with a horizontal dashed line at the C0 baseline. Alternatively, a scatter
% plot with Val WER on the x-axis and Test WER on the y-axis, with each
% condition as a labelled point.

\begin{figure}[H]
    \centering
    \includegraphics[width=0.95\linewidth]{placeholder_integrated_summary.png}
    \caption{Summary of all confirmatory conditions under the speaker-disjoint
    v2 protocol. Conditions below the horizontal dashed line (C0 baseline)
    represent improvements; conditions above represent degradations. SpecAugment
    (S2) is the only condition to achieve simultaneous improvement on both
    validation and test partitions.}
    \label{fig:integrated-summary}
\end{figure}


% =====================================================================
\subsection{Limitations}
\label{sec:results-limitations}
% =====================================================================

The following limitations apply to the results reported in this chapter and
should be considered when interpreting the findings:

\begin{enumerate}
    \item Most RQ1 curriculum conditions were screened with a single seed (seed
    42), and the seed-42 confidence interval family was not adjusted for
    multiplicity. The intervals therefore provide individual rather than
    family-wise error control.
    
    \item C3 is deterministic under zero dropout and strict ordering. Its
    apparent zero standard deviation across seeds is a structural artefact of
    checkpoint reuse, not evidence of zero training variance or three independent
    replications.
    
    \item The WAXAL test partition contributed to the selection of C3 for the
    A-series factorial. C3's test WER is therefore selection-stage evidence
    rather than an independent confirmation of its advantage.
    
    \item FLEURS bootstrap intervals cluster by source sentence prompt because
    speaker identities are unavailable. This clustering cannot account for
    unknown speaker-level dependence structures within the FLEURS corpus.
    
    \item The exploratory model-capacity and early SpecAugment experiments used
    the speaker-overlapping protocol and are not directly comparable with the
    confirmatory speaker-disjoint results.
    
    \item The A-series curriculum--augmentation interaction
    (difference-in-differences) was not formally cluster-bootstrapped as an
    interaction contrast; the reported interaction estimate is descriptive.
    
    \item All results are conditional on the specific Whisper Base model
    revision, Shona training data, three-epoch budget, speaker-disjoint v2
    split, preprocessing contract, and hyperparameter configuration used in this
    study. Generalisation to other model sizes, languages, training budgets, or
    evaluation protocols requires independent verification.
    
    \item The SpecAugment results (S-series) require verification of the S3, S4,
    and S5 cells from the trainer-state artifacts before final submission.
    Values reported for S2 and S4 in this draft should be confirmed as distinct
    or reconciled.
\end{enumerate}


\end{document}
